\documentclass[a4paper,12pt]{article}

% --- Use XeLaTeX font support ---

\usepackage{fontspec}
\setmainfont{TeX Gyre Termes} % nice Times-like font, installed by texlive-core
\setsansfont{TeX Gyre Heros} % optional sans-serif
\setmonofont{Fira Code}       % optional monospaced font
% --- Math packages ---
\usepackage{amsmath, amssymb, amsthm}
\usepackage{mathtools}

% --- Page layout ---
\usepackage[margin=1in]{geometry}

% --- Better lists ---
\usepackage{enumitem}
% --- Hyperlinks for references ---
\usepackage[colorlinks=true, linkcolor=blue, urlcolor=blue]{hyperref}

% Header: fancy layout like the rovided image
\usepackage{fancyhdr}
\setlength{\headheight}{26pt} % increase if fancyhdr warns
\pagestyle{fancy}
\fancyhf{} % clear header/footer
\rhead{\begin{tabular}{r} Mt.: 3052737\end{tabular}}
\chead{\begin{tabular}{c}Lösungen EA 2\end{tabular}}
\lhead{\begin{tabular}{l}61112 Lineare Algebra SS26 \end{tabular}}
\renewcommand{\headrulewidth}{0.8pt}
\fancyfoot[C]{- \thepage\ -}
% Ensure the plain page style (used by \maketitle) uses the same header

% --- Line spacing ---
\usepackage{setspace} % provides \doublespacing and \onehalfspacing

% --- Optional solutions environment ---
\newenvironment{solution}{
  \par\noindent\textbf{Solution:}\quad
}{\hfill$\blacksquare$\par\vspace{1em}}

% % --- Title info ---
% \title{Aufgabe \#1}
% \author{Your Name}
% \date{\today}
% 
\begin{document}
\raggedright
% Enable double (2.0) line spacing for the document body
\doublespacing
\noindent\textbf{Aufgabe 2.2} \\


\noindent\textbf{(a)} Die Matrix wird durch die Bilder der 
Basisvektoren bestimmt. 
Wir setzen die Standardbasisvektoren in $f$ ein:
$$
f(e_1) = f\begin{pmatrix} 1 \\ 0 \end{pmatrix} = \begin{pmatrix} 1 + 0 \\ (1+i)\cdot 1 - 0 \end{pmatrix} = \begin{pmatrix} 1 \\ 1+i \end{pmatrix}
$$
$$
f(e_2) = f\begin{pmatrix} 0 \\ 1 \end{pmatrix} = \begin{pmatrix} 0 + 1 \\ (1+i)\cdot 0 - 1 \end{pmatrix} = \begin{pmatrix} 1 \\ -1 \end{pmatrix}
$$
Die Bildvektoren sind direkt die Spalten der Matrix:
$${}_{\mathcal{E}}M_{\mathcal{E}}(f)=\begin{pmatrix} 1 & 1 \\ 1 + i & -1   \end{pmatrix}.$$

\bigskip
\noindent\textbf{(b)}
 Wenn man eine Darstellungsmatrix hat, dann ist $f$ genau ein Isomorphismus, wenn die Matrix invertierbar ist. Man erkennt sofot, dass die Matrix aus (a)  den vollen Rang hat (Rang 2), ist damit invertierbar, und f ist ein Isomorphismus.


\bigskip
\noindent\textbf{(c)} Die Transformationsformel lautet:
$${}_{\mathcal{B}}M_{\mathcal{B}}(f)={}_{\mathcal{B}}M_{\mathcal{E}}(id){}_{\mathcal{E}}M_{\mathcal{E}}(f){}_{\mathcal{E}}M_{\mathcal{B}}(id)$$

Die Spalten von ${}_{\mathcal{E}}M_{\mathcal{B}}(id)$ sind einfach die Vektoren der Basis \\
$$ {}_{\mathcal{E}}M_{\mathcal{B}}(id)= \begin{pmatrix} 0 & 1 \\ i & 2i \end{pmatrix}$$
Durch berechnen der Inverse von ${}_{\mathcal{E}}M_{\mathcal{B}}(id)$ erhalten wir ${}_{\mathcal{B}}M_{\mathcal{E}}(id)$:
$${}_{\mathcal{B}}M_{\mathcal{E}}(id) = \frac{1}{-i} \begin{pmatrix} 2i & -1 \\ -i & 0 \end{pmatrix} = i \begin{pmatrix} 2i & -1 \\ -i & 0 \end{pmatrix} = \begin{pmatrix} -2 & -i \\ 1 & 0 \end{pmatrix}$$
Zusammen: \\
$$M_\mathcal{B}^\mathcal{B}(f) = \begin{pmatrix} -2 & -i \\ 1 & 0 \end{pmatrix} \cdot \begin{pmatrix} 1 & 1 \\ 1+i & -1 \end{pmatrix} \cdot \begin{pmatrix} 0 & 1 \\ i & 2i \end{pmatrix}$$
Ausmultiplizieren und vereinfachen ergibt:
$$ M_\mathcal{B}^\mathcal{B}(f)= \begin{pmatrix} -1 - 2i & -3 - 5i \\ i & 1 + 2i \end{pmatrix}$$


\bigskip
\noindent\textbf{(d)} 
Gesucht sind Linearformen $b_i^*: \mathbb{C}^2 \to \mathbb{C}$, sodass $b_i^*(b_j) = \delta_{ij}$ gilt. Wie in Beispiel (2.4.14) kann man die Dualbasis aus den Zeilen der Inverse der  Matrix 
$A = \begin{pmatrix} 0 & 1 \\ i & 2i \end{pmatrix}$ ablesen.
$$A^{-1} = \begin{pmatrix} -2 & -i \\ 1 & 0 \end{pmatrix}$$ 

Die Dualbasisvektoren sind definiert durch:
$$b_1^*\begin{pmatrix} x \\ y \end{pmatrix} = -2x - iy$$
$$b_2^*\begin{pmatrix} x \\ y \end{pmatrix} = x$$
\end{document}
